\chapter{Conclusion et perspectives}
Un pare-feu entièrement fonctionnel a été déployé avec succès sur un serveur Ubuntu en utilisant \texttt{iptables}. Le pare-feu :
\begin{itemize}
    \item Bloque tout le trafic par défaut (approche de la liste blanche)
    \item N'autorise que SSH, HTTP et HTTPS
    \item Protège SSH contre les attaques par force brute en utilisant la limitation du débit de connexion
    \item Défend contre les attaques par inondation SYN, les paquets fragmentés et les paquets TCP invalides
    \item Enregistre les paquets suspects rejetés pour la surveillance
    \item Persiste après les redémarrages
\end{itemize}
Les scans nmap avant/après ont confirmé que le pare-feu fonctionne : 1021 ports sont passés de "fermés" à "filtrés", et le temps de scan a augmenté de 0.22 à 4.79 secondes car les paquets rejetés ne produisent aucune réponse et forcent nmap à attendre les délais d'attente.

\section{Perspectives}
\lipsum[11-12]
